\chapter{From reference solver to research experiments}
\label{chap:extensions}

The current solver is intentionally small.  Its purpose is not to compete with
high-performance many-body packages, but to provide a transparent numerical
contract for new representations.

\section{A safe extension sequence}

The following order preserves independent validation at each stage.
\begin{enumerate}
\item \textbf{Arbitrary one-body graphs.}  Already supported through
      Chapter~\ref{chap:generalhopping}; use this for longer-range hopping, disorder,
      and flux.
\item \textbf{Natural-orbital diagnostics.}  Diagonalize the one-body density
      matrix from Eq.~\eqref{eq:obdm}.  At first, treat its eigenvectors as an
      analysis tool rather than silently changing the interacting Hamiltonian.
\item \textbf{General two-body tensors.}  Implement Eq.~\eqref{eq:rotatedinteraction}
      with an explicit index and operator-ordering convention.  Check unitary
      basis invariance of the complete small-system spectrum.
\item \textbf{Wavelet and learned rotations.}  Transform both one- and two-body
      terms.  Compare energies, residuals, reduced density matrices, and
      truncation errors against the site-basis result.
\item \textbf{Entanglement diagnostics.}  Add subsystem maps and reduced density
      matrices only after fixing the tensor-product ordering used for a cut.
\item \textbf{External methods.}  Compare DMRG, neural-network, and transformer
      wavefunctions on identical finite chains and boundary conventions.
\end{enumerate}

\section{A reproducible experiment contract}

Every comparison should record
\begin{itemize}
\item $L$, $(N_\uparrow,N_\downarrow)$, $t$, $U$, and boundary condition;
\item the spin-orbital and basis enumeration conventions;
\item Hilbert-space dimension and whether the action is CSR or matrix-free;
\item solver tolerance, random seed, eigenvalue count, and measured residuals;
\item all orbital transforms and whether the interaction was transformed;
\item observables with their normalization conventions;
\item agreement with at least one analytic limit and the site-basis reference.
\end{itemize}
Energy alone is not a sufficient comparison.  Two states with similar energy
can have different correlations, symmetry content, or reduced density
matrices.

\section{Performance interpretation}

The half-filled dimensions in Chapter~\ref{chap:basis} explain the practical
range.  On the reference host, $L=8$ is quick and $L=10$ is an upper practical
target; $L=12$ is generally too slow and memory-heavy in pure Python.  A dense
general hopping matrix also has up to $L^2$ terms per spin, so a matrix-free
product can be much slower than the nearest-neighbor action at the same
$\basisdim$.

Optimization should follow profiling and retain a slow trusted path.  Useful
future optimizations include cached transition tables, compiled bit kernels,
symmetry-sector reductions, and parallel observable evaluation.  Each optimized
path should continue to pass the random-vector and full-spectrum comparisons
against the current implementation.

\section{Suggested exercises}

\begin{enumerate}
\item Add a site-dependent diagonal potential through $h_{ii}$ and verify the
      $U=0$ full spectrum from the one-particle eigenvalues.
\item Add a next-nearest-neighbor edge whose source and destination have an
      occupied site between them; predict and test the sign.
\item Compute natural-orbital occupations for $U/t=0,4,16$ on the open $L=6$
      chain.  Check their sum and spin symmetry.
\item Construct a random unitary $R$ at $U=0$ and verify that $h$ and
      $R^\dagger hR$ give identical many-body spectra.
\item Implement the four-index interaction for $L=2$ and verify full spectral
      invariance under a random orbital rotation.
\end{enumerate}

The last exercise is the natural bridge from this tutorial to serious orbital
representation research: it tests not only code, but whether the transformed
Hamiltonian still represents the same physics.
